\IfFileExists{stacks-project.cls}{%
\documentclass{stacks-project}
}{%
\documentclass{amsart}
}

% For dealing with references we use the comment environment
\usepackage{verbatim}
\newenvironment{reference}{\comment}{\endcomment}
%\newenvironment{reference}{}{}
\newenvironment{slogan}{\comment}{\endcomment}
\newenvironment{history}{\comment}{\endcomment}

% For commutative diagrams we use Xy-pic
\usepackage[all]{xy}

% We use 2cell for 2-commutative diagrams.
\xyoption{2cell}
\UseAllTwocells

% We use multicol for the list of chapters between chapters
\usepackage{multicol}

% This is generally recommended for better output
\usepackage{lmodern}
\usepackage[T1]{fontenc}

% For cross-file-references
\usepackage{xr-hyper}

% Package for hypertext links:
\usepackage{hyperref}

% For any local file, say "hello.tex" you want to link to please
% use \externaldocument[hello-]{hello}
\externaldocument[introduction-]{introduction}
\externaldocument[conventions-]{conventions}
\externaldocument[sets-]{sets}
\externaldocument[categories-]{categories}
\externaldocument[topology-]{topology}
\externaldocument[sheaves-]{sheaves}
\externaldocument[sites-]{sites}
\externaldocument[stacks-]{stacks}
\externaldocument[fields-]{fields}
\externaldocument[algebra-]{algebra}
\externaldocument[brauer-]{brauer}
\externaldocument[homology-]{homology}
\externaldocument[derived-]{derived}
\externaldocument[simplicial-]{simplicial}
\externaldocument[more-algebra-]{more-algebra}
\externaldocument[smoothing-]{smoothing}
\externaldocument[modules-]{modules}
\externaldocument[sites-modules-]{sites-modules}
\externaldocument[injectives-]{injectives}
\externaldocument[cohomology-]{cohomology}
\externaldocument[sites-cohomology-]{sites-cohomology}
\externaldocument[dga-]{dga}
\externaldocument[dpa-]{dpa}
\externaldocument[sdga-]{sdga}
\externaldocument[hypercovering-]{hypercovering}
\externaldocument[schemes-]{schemes}
\externaldocument[constructions-]{constructions}
\externaldocument[properties-]{properties}
\externaldocument[morphisms-]{morphisms}
\externaldocument[coherent-]{coherent}
\externaldocument[divisors-]{divisors}
\externaldocument[limits-]{limits}
\externaldocument[varieties-]{varieties}
\externaldocument[topologies-]{topologies}
\externaldocument[descent-]{descent}
\externaldocument[perfect-]{perfect}
\externaldocument[more-morphisms-]{more-morphisms}
\externaldocument[flat-]{flat}
\externaldocument[groupoids-]{groupoids}
\externaldocument[more-groupoids-]{more-groupoids}
\externaldocument[etale-]{etale}
\externaldocument[chow-]{chow}
\externaldocument[intersection-]{intersection}
\externaldocument[pic-]{pic}
\externaldocument[weil-]{weil}
\externaldocument[adequate-]{adequate}
\externaldocument[dualizing-]{dualizing}
\externaldocument[duality-]{duality}
\externaldocument[discriminant-]{discriminant}
\externaldocument[derham-]{derham}
\externaldocument[local-cohomology-]{local-cohomology}
\externaldocument[algebraization-]{algebraization}
\externaldocument[curves-]{curves}
\externaldocument[resolve-]{resolve}
\externaldocument[models-]{models}
\externaldocument[functors-]{functors}
\externaldocument[equiv-]{equiv}
\externaldocument[pione-]{pione}
\externaldocument[etale-cohomology-]{etale-cohomology}
\externaldocument[proetale-]{proetale}
\externaldocument[relative-cycles-]{relative-cycles}
\externaldocument[more-etale-]{more-etale}
\externaldocument[trace-]{trace}
\externaldocument[crystalline-]{crystalline}
\externaldocument[spaces-]{spaces}
\externaldocument[spaces-properties-]{spaces-properties}
\externaldocument[spaces-morphisms-]{spaces-morphisms}
\externaldocument[decent-spaces-]{decent-spaces}
\externaldocument[spaces-cohomology-]{spaces-cohomology}
\externaldocument[spaces-limits-]{spaces-limits}
\externaldocument[spaces-divisors-]{spaces-divisors}
\externaldocument[spaces-over-fields-]{spaces-over-fields}
\externaldocument[spaces-topologies-]{spaces-topologies}
\externaldocument[spaces-descent-]{spaces-descent}
\externaldocument[spaces-perfect-]{spaces-perfect}
\externaldocument[spaces-more-morphisms-]{spaces-more-morphisms}
\externaldocument[spaces-flat-]{spaces-flat}
\externaldocument[spaces-groupoids-]{spaces-groupoids}
\externaldocument[spaces-more-groupoids-]{spaces-more-groupoids}
\externaldocument[bootstrap-]{bootstrap}
\externaldocument[spaces-pushouts-]{spaces-pushouts}
\externaldocument[spaces-chow-]{spaces-chow}
\externaldocument[groupoids-quotients-]{groupoids-quotients}
\externaldocument[spaces-more-cohomology-]{spaces-more-cohomology}
\externaldocument[spaces-simplicial-]{spaces-simplicial}
\externaldocument[spaces-duality-]{spaces-duality}
\externaldocument[formal-spaces-]{formal-spaces}
\externaldocument[restricted-]{restricted}
\externaldocument[spaces-resolve-]{spaces-resolve}
\externaldocument[formal-defos-]{formal-defos}
\externaldocument[defos-]{defos}
\externaldocument[cotangent-]{cotangent}
\externaldocument[examples-defos-]{examples-defos}
\externaldocument[algebraic-]{algebraic}
\externaldocument[examples-stacks-]{examples-stacks}
\externaldocument[stacks-sheaves-]{stacks-sheaves}
\externaldocument[criteria-]{criteria}
\externaldocument[artin-]{artin}
\externaldocument[quot-]{quot}
\externaldocument[stacks-properties-]{stacks-properties}
\externaldocument[stacks-morphisms-]{stacks-morphisms}
\externaldocument[stacks-limits-]{stacks-limits}
\externaldocument[stacks-cohomology-]{stacks-cohomology}
\externaldocument[stacks-perfect-]{stacks-perfect}
\externaldocument[stacks-introduction-]{stacks-introduction}
\externaldocument[stacks-more-morphisms-]{stacks-more-morphisms}
\externaldocument[stacks-geometry-]{stacks-geometry}
\externaldocument[moduli-]{moduli}
\externaldocument[moduli-curves-]{moduli-curves}
\externaldocument[examples-]{examples}
\externaldocument[exercises-]{exercises}
\externaldocument[guide-]{guide}
\externaldocument[desirables-]{desirables}
\externaldocument[coding-]{coding}
\externaldocument[obsolete-]{obsolete}
\externaldocument[fdl-]{fdl}
\externaldocument[index-]{index}

% Theorem environments.
%
\theoremstyle{plain}
\newtheorem{theorem}[subsection]{Theorem}
\newtheorem{proposition}[subsection]{Proposition}
\newtheorem{lemma}[subsection]{Lemma}

\theoremstyle{definition}
\newtheorem{definition}[subsection]{Definition}
\newtheorem{example}[subsection]{Example}
\newtheorem{exercise}[subsection]{Exercise}
\newtheorem{situation}[subsection]{Situation}

\theoremstyle{remark}
\newtheorem{remark}[subsection]{Remark}
\newtheorem{remarks}[subsection]{Remarks}

\numberwithin{equation}{subsection}

% Macros
%
\def\lim{\mathop{\mathrm{lim}}\nolimits}
\def\colim{\mathop{\mathrm{colim}}\nolimits}
\def\Spec{\mathop{\mathrm{Spec}}}
\def\Hom{\mathop{\mathrm{Hom}}\nolimits}
\def\Ext{\mathop{\mathrm{Ext}}\nolimits}
\def\SheafHom{\mathop{\mathcal{H}\!\mathit{om}}\nolimits}
\def\SheafExt{\mathop{\mathcal{E}\!\mathit{xt}}\nolimits}
\def\Sch{\mathit{Sch}}
\def\Mor{\mathop{\mathrm{Mor}}\nolimits}
\def\Ob{\mathop{\mathrm{Ob}}\nolimits}
\def\Sh{\mathop{\mathit{Sh}}\nolimits}
\def\NL{\mathop{N\!L}\nolimits}
\def\CH{\mathop{\mathrm{CH}}\nolimits}
\def\proetale{{pro\text{-}\acute{e}tale}}
\def\etale{{\acute{e}tale}}
\def\QCoh{\mathit{QCoh}}
\def\Ker{\mathop{\mathrm{Ker}}}
\def\Im{\mathop{\mathrm{Im}}}
\def\Coker{\mathop{\mathrm{Coker}}}
\def\Coim{\mathop{\mathrm{Coim}}}

% Boxtimes
%
\DeclareMathSymbol{\boxtimes}{\mathbin}{AMSa}{"02}

%
% Macros for moduli stacks/spaces
%
\def\QCohstack{\mathcal{QC}\!\mathit{oh}}
\def\Cohstack{\mathcal{C}\!\mathit{oh}}
\def\Spacesstack{\mathcal{S}\!\mathit{paces}}
\def\Quotfunctor{\mathrm{Quot}}
\def\Hilbfunctor{\mathrm{Hilb}}
\def\Curvesstack{\mathcal{C}\!\mathit{urves}}
\def\Polarizedstack{\mathcal{P}\!\mathit{olarized}}
\def\Complexesstack{\mathcal{C}\!\mathit{omplexes}}
% \Pic is the operator that assigns to X its picard group, usage \Pic(X)
% \Picardstack_{X/B} denotes the Picard stack of X over B
% \Picardfunctor_{X/B} denotes the Picard functor of X over B
\def\Pic{\mathop{\mathrm{Pic}}\nolimits}
\def\Picardstack{\mathcal{P}\!\mathit{ic}}
\def\Picardfunctor{\mathrm{Pic}}
\def\Deformationcategory{\mathcal{D}\!\mathit{ef}}
\usepackage{cleveref}
\crefname{originallemma}{original lemma}{original lemmas}
\Crefname{originallemma}{Original lemma}{Original lemmas}
\makeatletter
\newlabel{algebra-lemma-ML-exact-sequence}{{0598}{source}{}{sourcepointer-1}{}}
\newlabel{algebra-lemma-ML-exact-sequence@cref}{{[originallemma][0][]0598}{[1][0][]source}{}{}{}}
\makeatother
\crefname{originallemma}{original lemma}{original lemmas}
\Crefname{originallemma}{Original lemma}{Original lemmas}
\makeatletter
\newlabel{algebra-lemma-compose-finite-type}{{00F4}{source}{}{sourcepointer-2}{}}
\newlabel{algebra-lemma-compose-finite-type@cref}{{[originallemma][0][]00F4}{[1][0][]source}{}{}{}}
\makeatother
\crefname{originallemma}{original lemma}{original lemmas}
\Crefname{originallemma}{Original lemma}{Original lemmas}
\makeatletter
\newlabel{algebra-lemma-surjective-mod-locally-nilpotent}{{07RD}{source}{}{sourcepointer-3}{}}
\newlabel{algebra-lemma-surjective-mod-locally-nilpotent@cref}{{[originallemma][0][]07RD}{[1][0][]source}{}{}{}}
\makeatother
\crefname{originallemma}{original lemma}{original lemmas}
\Crefname{originallemma}{Original lemma}{Original lemmas}
\makeatletter
\newlabel{descent-lemma-finite-type-local-source-fppf-algebra}{{0367}{source}{}{sourcepointer-4}{}}
\newlabel{descent-lemma-finite-type-local-source-fppf-algebra@cref}{{[originallemma][0][]0367}{[1][0][]source}{}{}{}}
\makeatother
\crefname{originalequation}{original equation}{original equations}
\Crefname{originalequation}{Original equation}{Original equations}
\makeatletter
\newlabel{equation-localize}{{0ANC}{source}{}{sourcepointer-5}{}}
\newlabel{equation-localize@cref}{{[originalequation][0][]0ANC}{[1][0][]source}{}{}{}}
\makeatother
\crefname{originaldefinition}{original definition}{original definitions}
\Crefname{originaldefinition}{Original definition}{Original definitions}
\makeatletter
\newlabel{example-finite-type-from-finite-type-ring-map}{{0ANM}{source}{}{sourcepointer-6}{}}
\newlabel{example-finite-type-from-finite-type-ring-map@cref}{{[originaldefinition][0][]0ANM}{[1][0][]source}{}{}{}}
\makeatother
\crefname{originalreference}{original reference}{original references}
\Crefname{originalreference}{Original reference}{Original references}
\makeatletter
\newlabel{item-axiom-1}{{0AND}{source}{}{sourcepointer-7}{}}
\newlabel{item-axiom-1@cref}{{[originalreference][0][]0AND}{[1][0][]source}{}{}{}}
\makeatother
\crefname{originalreference}{original reference}{original references}
\Crefname{originalreference}{Original reference}{Original references}
\makeatletter
\newlabel{item-axiom-2}{{0ANE}{source}{}{sourcepointer-8}{}}
\newlabel{item-axiom-2@cref}{{[originalreference][0][]0ANE}{[1][0][]source}{}{}{}}
\makeatother
\crefname{originalreference}{original reference}{original references}
\Crefname{originalreference}{Original reference}{Original references}
\makeatletter
\newlabel{item-axiom-3}{{0ANF}{source}{}{sourcepointer-9}{}}
\newlabel{item-axiom-3@cref}{{[originalreference][0][]0ANF}{[1][0][]source}{}{}{}}
\makeatother
\crefname{originallemma}{original lemma}{original lemmas}
\Crefname{originallemma}{Original lemma}{Original lemmas}
\makeatletter
\newlabel{lemma-representable-property-rings}{{0ANJ}{source}{}{sourcepointer-10}{}}
\newlabel{lemma-representable-property-rings@cref}{{[originallemma][0][]0ANJ}{[1][0][]source}{}{}{}}
\makeatother
\crefname{originalremark}{original remark}{original remarks}
\Crefname{originalremark}{Original remark}{Original remarks}
\makeatletter
\newlabel{remark-universal-property}{{0AJM}{source}{}{sourcepointer-11}{}}
\newlabel{remark-universal-property@cref}{{[originalremark][0][]0AJM}{[1][0][]source}{}{}{}}
\makeatother
\crefname{originalsection}{original section}{original sections}
\Crefname{originalsection}{Original section}{Original sections}
\makeatletter
\newlabel{section-morphisms-rings}{{0ANA}{source}{}{sourcepointer-12}{}}
\newlabel{section-morphisms-rings@cref}{{[originalsection][0][]0ANA}{[1][0][]source}{}{}{}}
\makeatother
\crefname{originallemma}{original lemma}{original lemmas}
\Crefname{originallemma}{Original lemma}{Original lemmas}
\makeatletter
\newlabel{lemma-completion-in-sub}{{0ANB}{source}{}{sourcepointer-13}{}}
\newlabel{lemma-completion-in-sub@cref}{{[originallemma][0][]0ANB}{[1][0][]source}{}{}{}}
\makeatother
\crefname{originaldefinition}{original definition}{original definitions}
\Crefname{originaldefinition}{Original definition}{Original definitions}
\makeatletter
\newlabel{situation-local-property}{{0CBA}{source}{}{sourcepointer-14}{}}
\newlabel{situation-local-property@cref}{{[originaldefinition][0][]0CBA}{[1][0][]source}{}{}{}}
\makeatother
\title{Topologically of finite type over a complete Tate ring: equivalent formulations and locality}
\author{The Stacks Project Authors}
\date{stacks-project:formal-spaces.tex:local-snapshot-20261012; GNU FDL 1.2 or later}
\begin{document}
\maketitle
\noindent\textit{Editorial note.} All delivered mathematical text is the author's own, copied from the pinned Stacks Project chapter file formal-spaces.tex; no statement or proof step was rewritten, reordered or removed.\par\medskip
\noindent\textit{Editorial note.} The last sentence of the proof is the author's 'We omit the proof of Axiom (3).' It is reproduced verbatim and is recorded as an obligation of the card: the equivalence of conditions (a) to (d) and locality axioms (1) and (2) are proved in full inside the counted span, and only the third locality axiom is left out.\par\medskip
\noindent\textit{Editorial note.} The statement refers to the cofinal-system formulation without recalling the definitions of tautness, weak ideal of definition or topologically of finite type; those are the chapter's standing notions and the same-chapter situation governing the locality axioms is delivered as a context span, so the reader sees the shape of the diagrams the second half of the proof works with.\par\medskip
\noindent\textit{Editorial note.} The proof displays three commutative diagrams and two exact sequences; all lie inside the counted span and are reproduced verbatim.\par\medskip
\noindent\textit{Editorial note.} The counted spans carry 13 navigation stubs; no bibliography citation occurs in the spans.\par\medskip
\noindent\textit{Editorial note.} The preamble span points at the chapter's own preamble.tex shipped with the frozen source, used verbatim as staged.\par\medskip
\noindent\textit{Allowed standard prerequisites (editorial).}\begin{itemize}
\item The situation and the three locality axioms for a property of arrows of complete weakly admissible topological rings. Same chapter, Situation situation-local-property, invoked at proof lines 6623-6624, together with items item-axiom-1, item-axiom-2 and item-axiom-3 referred to in the statement. Fixes what has to be verified in the second half of the lemma.
\item Representation of an arrow as an inverse limit of its quotients. Same chapter, Lemma lemma-representable-property-rings, invoked at proof lines 6600-6601 and again at proof lines 6631-6633. Supplies the compatible system of quotients and the reduction that lets the two given arrows be assumed taut.
\item The universal property of the restricted power series algebra and surjectivity modulo locally nilpotent ideals. Same chapter, Remark remark-universal-property at proof lines 6610-6611, and Algebra, Lemma algebra-lemma-surjective-mod-locally-nilpotent at proof lines 6616-6618. Produces the continuous map from the restricted power series algebra onto A and the surjectivity of its reductions.
\item Exactness of an inverse limit of short exact sequences with surjective transition maps. Algebra, Lemma algebra-lemma-ML-exact-sequence, invoked at proof lines 6625-6627. Gives the final exactness which shows the constructed map is surjective and open.
\item Finiteness of a composition of ring maps, and descent of finite type along a faithfully flat etale ring map. Algebra, Lemma algebra-lemma-compose-finite-type and Descent, Lemma descent-lemma-finite-type-local-source-fppf-algebra, invoked at proof lines 6650-6655. Gives the two directions of locality axiom (1) and (2) for the quotients of the two arrows.
\end{itemize}\medskip
\ifdefined\hypertarget\hypertarget{sourcepointer-1}{}\fi
\noindent\textit{Source reference.} Original lemma 0598: sources/stacks-chow/tags.txt, source lines 6806--6806. Pointer outside the counted span, invoked once for context; the load-bearing argument of the counted statement is the author's proof reproduced in full in the proof span.\par\medskip
\ifdefined\hypertarget\hypertarget{sourcepointer-2}{}\fi
\noindent\textit{Source reference.} Original lemma 00F4: sources/stacks-chow/tags.txt, source lines 558--558. Pointer outside the counted span, invoked once for context; the load-bearing argument of the counted statement is the author's proof reproduced in full in the proof span.\par\medskip
\ifdefined\hypertarget\hypertarget{sourcepointer-3}{}\fi
\noindent\textit{Source reference.} Original lemma 07RD: sources/stacks-chow/tags.txt, source lines 9863--9863. Pointer outside the counted span, invoked once for context; the load-bearing argument of the counted statement is the author's proof reproduced in full in the proof span.\par\medskip
\ifdefined\hypertarget\hypertarget{sourcepointer-4}{}\fi
\noindent\textit{Source reference.} Original lemma 0367: sources/stacks-chow/tags.txt, source lines 4154--4154. Pointer outside the counted span, invoked once for context; the load-bearing argument of the counted statement is the author's proof reproduced in full in the proof span.\par\medskip
\ifdefined\hypertarget\hypertarget{sourcepointer-5}{}\fi
\noindent\textit{Source reference.} Original equation 0ANC: sources/stacks-formal-spaces/pinned-original-source.tex, source lines 5063--5063. Pointer outside the counted span, invoked once for context; the load-bearing argument of the counted statement is the author's proof reproduced in full in the proof span.\par\medskip
\ifdefined\hypertarget\hypertarget{sourcepointer-6}{}\fi
\noindent\textit{Source reference.} Original definition 0ANM: sources/stacks-formal-spaces/pinned-original-source.tex, source lines 5884--5884. Pointer outside the counted span, invoked once for context; the load-bearing argument of the counted statement is the author's proof reproduced in full in the proof span.\par\medskip
\ifdefined\hypertarget\hypertarget{sourcepointer-7}{}\fi
\noindent\textit{Source reference.} Original reference 0AND: sources/stacks-formal-spaces/pinned-original-source.tex, source lines 5086--5086. Pointer outside the counted span, invoked once for context; the load-bearing argument of the counted statement is the author's proof reproduced in full in the proof span.\par\medskip
\ifdefined\hypertarget\hypertarget{sourcepointer-8}{}\fi
\noindent\textit{Source reference.} Original reference 0ANE: sources/stacks-formal-spaces/pinned-original-source.tex, source lines 5090--5090. Pointer outside the counted span, invoked once for context; the load-bearing argument of the counted statement is the author's proof reproduced in full in the proof span.\par\medskip
\ifdefined\hypertarget\hypertarget{sourcepointer-9}{}\fi
\noindent\textit{Source reference.} Original reference 0ANF: sources/stacks-formal-spaces/pinned-original-source.tex, source lines 5095--5095. Pointer outside the counted span, invoked once for context; the load-bearing argument of the counted statement is the author's proof reproduced in full in the proof span.\par\medskip
\ifdefined\hypertarget\hypertarget{sourcepointer-10}{}\fi
\noindent\textit{Source reference.} Original lemma 0ANJ: sources/stacks-formal-spaces/pinned-original-source.tex, source lines 5466--5466. Pointer outside the counted span, invoked once for context; the load-bearing argument of the counted statement is the author's proof reproduced in full in the proof span.\par\medskip
\ifdefined\hypertarget\hypertarget{sourcepointer-11}{}\fi
\noindent\textit{Source reference.} Original remark 0AJM: sources/stacks-formal-spaces/pinned-original-source.tex, source lines 6287--6287. Pointer outside the counted span, invoked once for context; the load-bearing argument of the counted statement is the author's proof reproduced in full in the proof span.\par\medskip
\ifdefined\hypertarget\hypertarget{sourcepointer-12}{}\fi
\noindent\textit{Source reference.} Original section 0ANA: sources/stacks-formal-spaces/pinned-original-source.tex, source lines 4953--4953. Pointer outside the counted span, invoked once for context; the load-bearing argument of the counted statement is the author's proof reproduced in full in the proof span.\par\medskip
\ifdefined\hypertarget\hypertarget{sourcepointer-13}{}\fi
\noindent\textit{Source reference.} Original lemma 0ANB: sources/stacks-formal-spaces/pinned-original-source.tex, source lines 5029--5029. Pointer outside the counted span, invoked once for context; the load-bearing argument of the counted statement is the author's proof reproduced in full in the proof span.\par\medskip
\ifdefined\hypertarget\hypertarget{sourcepointer-14}{}\fi
\noindent\textit{Source reference.} Original definition 0CBA: sources/stacks-formal-spaces/pinned-original-source.tex, source lines 5059--5059. Pointer outside the counted span, invoked only from a delivered context span as a forward reference; the load-bearing argument of the counted statement is the author's proof reproduced in full in the proof span.\par\medskip
\begin{lemma}
\label{lemma-quotient-restricted-power-series}
Let $B \to A$ be an arrow of $\textit{WAdm}^{count}$, see
Section \ref{section-morphisms-rings}.
The following are equivalent
\begin{enumerate}
\item[(a)] $B \to A$ is taut and $B/J \to A/I$ is of finite type for
every weak ideal of definition $J \subset B$ where $I \subset A$ is the
closure of $JA$,
\item[(b)] $B \to A$ is taut and $B/J_\lambda \to A/I_\lambda$
is of finite type for a cofinal system $(J_\lambda)$
of weak ideals of definition of $B$ where
$I_\lambda \subset A$ is the closure of $J_\lambda A$,
\item[(c)] $B \to A$ is taut and $A$ is topologically of finite
type over $B$,
\item[(d)] $A$ is isomorphic as a topological $B$-algebra to a quotient of
$B\{x_1, \ldots, x_n\}$ by a closed ideal.
\end{enumerate}
Moreover, these equivalent conditions define a local property,
i.e., they satisfy
Axioms (\ref{item-axiom-1}), (\ref{item-axiom-2}), (\ref{item-axiom-3}).
\end{lemma}

\begin{proof}
The implications (a) $\Rightarrow$ (b), (c) $\Rightarrow$ (a),
(d) $\Rightarrow$ (c) are straightforward from the definitions.
Assume (b) holds and let $J \subset B$ and $I \subset A$ be as in (a).
Choose a commutative diagram
$$
\xymatrix{
A \ar[r] & \ldots \ar[r] & A_3 \ar[r] & A_2 \ar[r] & A_1 \\
B \ar[r] \ar[u] & \ldots \ar[r] & B/J_3 \ar[r] \ar[u] &
B/J_2 \ar[r] \ar[u] & B/J_1 \ar[u]
}
$$
such that $A_{n + 1}/J_nA_{n + 1} = A_n$ and such that $A = \lim A_n$ as in
Lemma \ref{lemma-representable-property-rings}.
For every $m$ there exists a $\lambda$ such that $J_\lambda \subset J_m$.
Since $B/J_\lambda \to A/I_\lambda$ is of finite type, this implies
that $B/J_m \to A/I_m$ is of finite type.
Let $\alpha_1, \ldots, \alpha_n \in A_1$ be generators of $A_1$ over
$B/J_1$. Since $A$ is a countable limit of a system with surjective
transition maps, we can find $a_1, \ldots, a_n \in A$ mapping to
$\alpha_1, \ldots, \alpha_n$ in $A_1$. By
Remark \ref{remark-universal-property} we find a continuous map
$B\{x_1, \ldots, x_n\} \to A$ mapping $x_i$ to $a_i$. This map
induces surjections $(B/J_m)[x_1, \ldots, x_n] \to A_m$ by
Algebra, Lemma \ref{algebra-lemma-surjective-mod-locally-nilpotent}.
For $m \geq 1$ we obtain a short exact sequence
$$
0 \to K_m \to (B/J_m)[x_1, \ldots, x_n] \to A_m \to 0
$$
The induced transition maps $K_{m + 1} \to K_m$ are surjective because
$A_{m + 1}/J_mA_{m + 1} = A_m$. Hence the inverse limit of these
short exact sequences is exact, see
Algebra, Lemma \ref{algebra-lemma-ML-exact-sequence}.
Since $B\{x_1, \ldots, x_n\} = \lim (B/J_m)[x_1, \ldots, x_n]$
and $A = \lim A_m$
we conclude that $B\{x_1, \ldots, x_n\} \to A$ is surjective and open.
As $A$ is complete the kernel is a closed ideal. In this way we see that
(a), (b), (c), and (d) are equivalent.

\medskip\noindent
Let a diagram (\ref{equation-localize}) as in
Situation \ref{situation-local-property}
be given. By Example \ref{example-finite-type-from-finite-type-ring-map}
the maps $A \to (A')^\wedge$ and $B \to (B')^\wedge$
satisfy (a), (b), (c), and (d). Moreover, by
Lemma \ref{lemma-representable-property-rings}
in order to prove Axioms (\ref{item-axiom-1}) and (\ref{item-axiom-2})
we may assume both $B \to A$ and $(B')^\wedge \to (A')^\wedge$
are taut. Now pick a weak ideal of definition $J \subset B$. Let
$J' \subset (B')^\wedge$, $I \subset A$, $I' \subset (A')^\wedge$
be the closure of $J(B')^\wedge$, $JA$, $J(A')^\wedge$.
By what was said above, it suffices to consider the commutative
diagram
$$
\xymatrix{
A/I \ar[r] & (A')^\wedge/I' \\
B/J \ar[r] \ar[u]^{\overline{\varphi}} &
(B')^\wedge/J' \ar[u]_{\overline{\varphi}'}
}
$$
and to show (1) $\overline{\varphi}$ finite type
$\Rightarrow \overline{\varphi}'$
finite type, and (2) if $A \to A'$ is faithfully flat, then
$\overline{\varphi}'$ finite type $\Rightarrow \overline{\varphi}$
finite type. Note that $(B')^\wedge/J' = B'/JB'$ and
$(A')^\wedge/I' = A'/IA'$ by the construction of the topologies on
$(B')^\wedge$ and $(A')^\wedge$. In particular the horizontal
maps in the diagram are \'etale. Part (1) now follows from
Algebra, Lemma \ref{algebra-lemma-compose-finite-type}
and part (2) from
Descent, Lemma \ref{descent-lemma-finite-type-local-source-fppf-algebra}
as the ring map $A/I \to (A')^\wedge/I' = A'/IA'$ is faithfully flat
and \'etale.

\medskip\noindent
We omit the proof of Axiom (\ref{item-axiom-3}).
\end{proof}

\end{document}
